\section{Complete primitive-action consumers}\label{sec:primitive-consumers}
Throughout this section $Z_J(U)=\sum_{N\lhd U}|\Epi(J,U/N)|$.
Every quotient is literal, and $J$ is the same complete source throughout.

\begin{lemma}[One derived source for soluble affine targets]
\label{lem:derived-affine}
Let $U=V\rtimes R$, where $V=\F_p^d$ and $R\ne1$ is soluble,
faithful and irreducible. Put $\End_R(V)=\F_{p^f}$,
$c_R=|N_{\GL_d(p)}(R)|/|C_{\GL_d(p)}(R)|$, and
$\eta=f\log(p)/(pd)$. For every $J\leq S_b$,
\[
 Z_J(U)=Z_J(R)+|\Epi(J,U)|,
 \qquad |\Epi(J,U)|\leq|V|c_R2^{\eta b}.
\]
More precisely, if $t$ is the derived length of $R$ and $a_X$
are the multiplicities in the head of $H_1(J^{(t)},\F_p)$ of
the distinct pulled-back natural $R$-modules, then
\begin{equation}\label{eq:derived-affine}
 |\Epi(J,U)|=|V|c_R\sum_X\bigl(p^{fa_X}-1\bigr).
\end{equation}
\end{lemma}
\begin{proof}
Put $F=J^{(t)}$, $Q=J/F$. The nontrivial last derived group
$R^{(t-1)}$ is abelian and has order prime to $p$, because
$O_p(R)=1$ for a faithful irreducible linear group. It has no
fixed vector in $V$. For an onto $\theta:Q\to R$, the
prime-to-$p$ part of the abelian group $Q^{(t-1)}$ maps onto
$R^{(t-1)}$. Averaging and inflation--restriction therefore
give $H^i(Q,V_\theta)=0$ for all $i\geq0$.
Restriction identifies
$H^1(J,V_\theta)$ with
$\Hom_Q(H_1(F,\F_p),V_\theta)$, and each restriction has
exactly $|V|$ literal cocycles. A nonzero restriction is onto
$V$ by irreducibility, and is exactly the condition that the
lift to $U$ is onto: $U^{(t)}=V$.

Two top maps give isomorphic modules precisely when they differ
by an automorphism of $R$ induced by $N_{\GL_d(p)}(R)$.
Each such orbit has $c_R$ members, proving
\eqref{eq:derived-affine}. All types lie in one head, so
$d\sum_Xa_X\leq d_p(F)\leq b/p$ and
$\sum_X(p^{fa_X}-1)\leq p^{f\sum_Xa_X}-1$.
Finally, an original normal subgroup of $U$ either contains $V$
or intersects it trivially. In the latter case it centralizes
$V$; since $C_U(V)=V$, it is trivial. This proves the exact
full-menu identity and the bound.
\end{proof}

\begin{theorem}[All sufficiently small-order actual actions]
\label{thm:small-order-consumer}
All transitive nonbinary actions of degree $w\geq64$ with
$\log|U|\leq w/8$ have one complete ordinary forward kernel
\[
 K_{\rm so}(n,b)=\Delta(n,b)
 \sum_{\substack{[U]:\deg U=n-b\geq64\\\log|U|\leq(n-b)/8}}
 \frac{2^{(\log|U|)^2+(\log|U|)(b/2+1)}}{a(U)},
\]
with row at most $2\cdot2^{-n}$ eventually and zero scalar.
In particular this covers every primitive affine action of
degree at least $1024$.
\end{theorem}
\begin{proof}
The ordinary permutation-generator bound
\cite[Proposition 2.1]{HoltRoneyDougal2013} gives
$d(J)\leq b/2+1$. For $\ell=\log|U|$, the number of target
normal subgroups is at most $2^{\ell^2}$, by padding a generating
tuple of each subgroup to length $\lfloor\ell\rfloor$.
The images of source generators then prove the displayed fibre
envelope. Original pointing gives this kernel directly.

The uniform coefficient/Gaussian comparison gives
\[
 \log\Delta(n,b)\leq-\frac{(w-(w\bmod2))b}{8}
             -\frac{w^2}{16}+O(w\log(n+2)).
\]
The envelope costs at most $w^2/64+wb/16+w/8$.
The action-class entropy is $o(w^2)$ by the transitive
enumeration theorem. For $w\geq64$, the negative terms are
at most $-wn/32-w^2/64$; uniformly absorb class entropy into
$w^2/128+O(w)$ and the remaining terms into $wn/64$.
Hence $K_{\rm so}(n,n-w)\leq2^{-wn/64}$. Sum over all
$w\geq64$.
For affine degree $w=p^d$, $\log|U|<(d+1)\log w
\leq(\log w)^2+\log w\leq w/8$ when $w\geq1024$.
The last inequality follows at $\log w=10$ and thereafter
by differentiating twice. A primitive affine group of that
degree cannot be binary: in characteristic two an irreducible
binary linear group has only the trivial simple module.
\end{proof}

\begin{theorem}[The soluble primitive affine range]
\label{thm:soluble-primitive-consumer}
All soluble primitive affine actions of degrees $5\leq w<1024$,
except the natural action $C_5\rtimes C_4$ of degree five,
have a complete ordinary forward row $O(2^{-an})$ and a
complete scalar $O(2^{-a'n^2})$, for fixed $a,a'>0$.
\end{theorem}
\begin{proof}
For each actual $U=\F_p^d\rtimes R$, use
Lemma~\ref{lem:derived-affine}. A regular prime action has
$Z_J(C_p)=|\Hom(J,C_p)|\leq p^{b/p}$; for $p\geq5$,
$\log(p)/p<(p-1)/8$. Suppose $R\ne1$.
If $p=2$, $w\geq8$, use the faithful action of $R$ on the
$w-1$ nonzero vectors. It has degree $v=w-1<w$ and
$\eta\leq1/2<w/8$, so Theorem~\ref{thm:finite-comparator}
applies with comparator $R$.

For odd $p$ and abelian $R$, irreducibility makes $R$ cyclic
in the multiplicative group of a field. Thus
$Z_J(R)=|\Hom(J,R)|\leq|J/J'|\leq3^{b/3}$.
At $w\geq7$, its exponent and $\eta$ are smaller than
$(w-1)/8$. At degree five the complements are $1,C_2,C_4$.
For $C_2$ use its degree-two comparator; the $C_4$ case is
the specified exception.

At degree nine let $Z=R\cap\{\pm I\}$ and $P=R/Z$,
acting faithfully on the four projective points. For each
original normal subgroup of $R$, a fixed top map has at
most $|\Hom(J,C_2)|$ lifts through its central kernel.
If $h_2(J)=|\Hom(J,C_2)|$ and $e=d_2(P)$, then
\[
 h_2(J)Z_J(P)\leq2^eZ_J(P\times C_2).
\]
Indeed, for one quotient $P/N$ with $c=|\Hom(P/N,C_2)|$,
the two product kernels contribute
$[1+h_2(J)-c]|\Epi(J,P/N)|$, and
$h_2(J)\leq2^e[1+h_2(J)-c]$ whenever the epimorphism
count is nonzero. These kernels are distinct for each $N$.
Consequently
$Z_J(R)\leq\nu(R)2^{d_2(P)}Z_J(P\times C_2)$.
The comparator has faithful degree six, smaller than $h(9)=8$;
the affine exponent is at most $\log(3)/3<1$.

At degree twenty-five the scalar subgroup of $R\leq\GL_2(5)$
has order dividing four. Its projective image is a soluble
subgroup of $\operatorname{PGL}_2(5)\cong S_5$, of order at
most twenty-four. An intransitive subgroup has this bound
immediately; a soluble transitive subgroup of prime degree five
has a regular normal $C_5$ and lies in its order-twenty normalizer.
Each central quotient lift costs at most
$|\Hom(J,C_4)|\leq2^{b/2}$. Counting all fixed target normals
and using $d(J)\leq b/2+1$ gives
$Z_J(R)\leq D_R2^{(\log48)b/2}$.
Its exponent is strictly below $h(25)/8=3$.

At degree twenty-seven the complete irreducible-subgroup
classification of Roney--Dougal and Unger
\cite{RoneyDougalUnger2003Affine}, in the pinned PrimGrp ordering
\cite{PrimGrp344}, has nine soluble affine rows.  Their point-stabilizer
orders are
\[
  12,13,24,24,24,26,39,48,78.
\]
Thus every soluble irreducible $R\leq\GL_3(3)$ has order at most
seventy-eight.  This is a direct catalogue consequence; no additional
Clifford-theoretic assertion is used.
Hence $Z_J(R)\leq D_R2^{(\log78)b/2}$, with
$(\log78)/2<26/8$.

In every other odd prime-power degree, the base-size theorem
\cite[Theorem 1.1]{HalasiMaroti2016} gives
$|R|\leq w^2$ for $p\geq5$ and $|R|\leq w^3$ for $p=3$.
Its complete-reducibility and $p$-solubility hypotheses follow
from irreducibility and solubility here. The remaining ranges
start at $w=49$ and $81$. The inequalities
$49^8<2^{48}$ and $81^{12}<2^{80}$, and monotonicity of
$w-1-8\log w$ and $w-1-12\log w$, give
$\log|R|/2<(w-1)/8$ in both ranges. The fixed-target ordinary
generator envelope applies as before.

These cases exhaust the finite set of action classes. Every
pointwise exponent has a strict physical gap; every comparator
has $0<v<h(w)$ and $\eta<h(w)/8$. Their original kernels
are $\Delta(n,n-w)$ times their respective fibre envelopes
divided by $a(U)$. Apply Theorem~\ref{thm:finite-comparator}
and take finite positive minima of these gaps. The pointwise
classes have scalar zero. This proves complete row and scalar
rates without assuming any normalized count is bounded.
\end{proof}

\begin{proposition}[Natural symmetric-four and Frobenius-five actions]
\label{prop:s4-f20-consumers}
The natural $S_4$ and $C_5\rtimes C_4$ actions have complete
ordinary forward rows $O(2^{-an})$ and complete scalars
$O(2^{-a'n^2})$, for positive constants.
\end{proposition}
\begin{proof}
For $S_4=\F_2^2\rtimes S_3$,
Lemma~\ref{lem:derived-affine} gives
$Z_J(S_4)\leq Z_J(S_3)+24\cdot2^{b/4}$.
The comparator $S_3$ has faithful degree three, less than
the original width four, and $1/4<4/8$. The finite comparator
theorem applies, using the original normalizer $24$.

For the natural degree-five action put
$z(J)=Z_J(C_4)=|\Hom(J,C_4)|$ and $\eta=\log(5)/5$.
The same lemma gives $Z_J(U)\leq z(J)+5\cdot2^{\eta b}$.
Declare $z(J)\leq2^{3b/8}$ cold. Its exact kernel is
\[
 \frac{\Delta(n,n-5)}{20}
       \left(2^{3(n-5)/8}+5\cdot2^{\eta(n-5)}\right),
\]
whose gap is positive since $\eta<15/32<1/2$.
On the hot event, the whole original fibre is at most
$6\cdot2^{(\eta-3/8)b}z(J)$.
Use Theorem~\ref{thm:C4-moment} with
$r=\max(1,\lceil b/8\rceil)$. Its weighted first moment,
including the hot cutoff, has quadratic exponent
\[
 b^2/16+br/4+r^2/2-(r-1)3b/8+(\eta-3/8)b-n^2/16
 =-b^2/128+(r-b/8)^2/2+(\eta-5/8)b-25/16.
\]
The actual scalar includes exactly the original factor
$1/(20b!p_nG_{\lfloor n/2\rfloor})$ on the hot source sum.
Its logarithmic normalization costs $O(n\log(n+2))$;
the moment error is $o(n^2)$. Thus it is at most
$2^{-n^2/512}$ eventually. The hot event is empty at $b=0$.
\end{proof}

\begin{theorem}[The nonsoluble primitive affine range]
\label{thm:nonsoluble-primitive-consumer}
All primitive affine actions $U=\F_p^d\rtimes R$ of degrees
$5\leq w=p^d\leq1024$ with nonsoluble $R$ have one complete
pointwise ordinary forward row at most $2^{-n/8}$ eventually
and scalar zero.
\end{theorem}
\begin{proof}
By Theorem~\ref{thm:composition-fibre} and
Proposition~\ref{prop:irreducible-affine}, the complete fibre
has the bound
\[
 Z_J(U)\leq(1+A_U)C_R(1+b\log(b+1))^{D_R}
              2^{(\kappa(R)+\eta_U)b},
\]
where $A_U=|Z^1(R,\F_p^d)|$,
$\eta_U=f\log(p)/(dp)$, and $\F_{p^f}=\End_R(\F_p^d)$.
Because $R$ is nonabelian, $f$ is a proper divisor of $d$;
otherwise the faithful image would be scalar over its commuting
field. Thus $d\geq2$ and $\eta_U\leq\log(p)/(2p)$.
Write $\gamma=\log(3)/3<17/32$.

The primitive composition bound
\cite[Theorem 1.3]{GlasbyPraegerRosaVerret2018} gives
$c(U)\leq(8/3)\log w-4/3$. Since $c(U)=d+c(R)$ and
$R$ has a nonabelian composition factor, the number of abelian
composition factors of $R$ is at most
$(8/3)\log w-7/3-d$. Hence
\[
 \kappa(R)+\eta_U
 \leq\gamma\bigl((8/3)\log w-11/6-d\bigr)
 \leq\frac{17}{32}\bigl((8/3)\log w-23/6\bigr).
\]
For $x\geq64$, the difference between $(x-1)/8$ and the
last expression is increasing: using $\log_e2>1/2$, its
derivative exceeds $1/8-17/(6x)>0$. At $x=64$ it is
$271/192>1/4$. Thus
$h(w)/8-\kappa(R)-\eta_U>1/4$ throughout that range.

At $w=32$, the integer composition bound gives at most six
abelian factors in $R$, so the exponent is at most
$6(17/32)+1/4=55/16$, leaving gap $9/16$.
At $w=49$, use $7^{16}<2^{45}$ to obtain at most ten
abelian factors, giving exponent at most $357/64$ and
gap $27/64$.
At $w=25$, quotient the scalar subgroup of order dividing four.
The projective image is a nonsoluble subgroup of
$\operatorname{PGL}_2(5)\cong S_5$, hence $A_5$ or $S_5$.
For completeness, an intransitive subgroup of $S_5$ is soluble;
a transitive subgroup has one or six Sylow-five subgroups.
The former case lies in the order-twenty normalizer. In the latter,
its order is $30,60,$ or $120$; order thirty would give an
impossible faithful degree-four action of $S_5$ by cosets.
The other two orders give $A_5,S_5$. Therefore $R$ has at
most three abelian factors, all $C_2$, so its exponent is
at most $3/2+17/64=113/64<3-1/4$.

The only remaining prime-power degrees below sixty-four are
$8,16,27$: prime-degree stabilizers are cyclic, and
$\GL_2(2),\GL_2(3)$ are soluble. The complete primitive
action catalogue at these three degrees gives the following
finite composition-factor data for the literal stabilizers:
\begin{center}
\begin{tabular}{cclc}
\toprule
degree&primitive indices&abelian factors of $R$&bound on $\kappa(R)$\\
\midrule
8&3&none&0\\
16&11--20&at most one $C_2$ and one $C_3$&$33/32$\\
27&10,11&at most one $C_2$&$1/2$\\\bottomrule
\end{tabular}
\end{center}
The affine classification is that of
\cite{RoneyDougalUnger2003Affine}; the representative conventions and
cross-catalogue indexing are those of
\cite{PrimGrp344,RoneyDougal2005Primitive}. The displayed
composition series supply respectively exponent bounds
$1/4,41/32,49/64$, all with physical gap greater than $1/4$.
The catalogue completeness is used only at these three degrees;
the other ranges were covered analytically.

Finally use the original pointed kernel
\[
 \Delta(n,b)\sum_{\substack{[U]\ \mathrm{in\ this\ family}\\w=n-b}}
 \frac{(1+A_U)C_R(1+b\log(b+1))^{D_R}
                  2^{(\kappa(R)+\eta_U)b}}{a(U)}.
\]
The family is fixed and finite, so all target constants and
polynomial exponents are fixed. Its complete row is
$O(n^D)2^{-n/4}\leq2^{-n/8}$ eventually.
The full cocycle factor $A_U$ remains in this argument; it
is not set equal to $|\F_p^d|$ for nonsoluble stabilizers.
\end{proof}
